\documentclass[11pt]{article}
\usepackage{geometry}
\geometry{verbose}
\setlength{\parskip}{6pt}
\setlength{\parindent}{0pt}
\usepackage{color}
\usepackage{amsmath, amssymb}
\usepackage[authoryear]{natbib}
\usepackage[unicode=true,
 bookmarks=false,
 breaklinks=false,pdfborder={0 0 1},backref=section,colorlinks=false]
 {hyperref}

\makeatletter
\newcommand{\R}{\mathbb{R}}
\newcommand{\E}{\mathbb{E}}
\newcommand{\Cov}{\mathrm{Cov}}
\newcommand{\Var}{\mathrm{Var}}
\providecommand{\tabularnewline}{\\}
\usepackage{graphicx}
\usepackage{float}
\usepackage{url}
\usepackage{array}
\usepackage{enumitem}

\title{\Huge Assignment 7}
\author{\Large Chris Conlon}
\date{\Large Fall 2026 --- due at the start of Session 13}

\makeatother

\begin{document}

\title{Assignment 7}
\maketitle
\begin{center}
\begin{tabular*}{0.9\textwidth}{@{\extracolsep{\fill}}@{\extracolsep{\fill}}l@{\extracolsep{\fill}}l@{\extracolsep{\fill}}l}
Econometrics I & $\qquad$ & Professor Chris Conlon\tabularnewline
NYU Stern &  & Email: \href{mailto:cconlon@stern.nyu.edu}{cconlon@stern.nyu.edu}\tabularnewline
\end{tabular*}
\par\end{center}

\textbf{Instructions.} Two problems on binary choice and selection (Session 12). \texttt{R} with
\texttt{fixest} (\texttt{feglm} for logit and probit) or Python with \texttt{pyfixest}. Turn in a PDF
with your results and reasoning, and your code.

\subsection*{1 Mortgage denials in Boston}

The \texttt{Hmda} data (\texttt{Ecdat} package in R; also in the course repository) contain
$2{,}380$ mortgage applications from the 1990 Boston HMDA study (Munnell et al.\ 1996). The outcome is
\texttt{deny}. Use the specification from the Session 12 notebook: \texttt{black}, debt-to-income
(\texttt{dir}), housing-expense-to-income (\texttt{hir}), loan-to-value indicators for medium
($0.8 \le \texttt{lvr} \le 0.95$) and high ($> 0.95$), consumer and mortgage credit scores
(\texttt{ccs}, \texttt{mcs}), public bad credit record (\texttt{pbcr}), and denied mortgage insurance
(\texttt{dmi}).

\begin{enumerate}[label=(\alph*)]
\item Report the denial rate for Black and for white applicants. Then estimate the linear
probability model, the probit, and the logit with the covariates above. Present the three in one
table with robust standard errors for the LPM.
\item The logit coefficient on \texttt{black} is roughly $1.8$ times the probit coefficient. Why is that
the expected ratio, and why does it not mean the models disagree?
\item Compute the effect of \texttt{black} on the denial probability three ways: the LPM coefficient;
the logit \emph{average} marginal effect (average over applicants of the predicted probability with
\texttt{black} set to $1$ minus set to $0$); and the logit effect \emph{at the means} of the
covariates. Using the shape of the logistic CDF, explain when the last two differ (where would the
``average applicant'' have to sit for the effect at the means to be much smaller than the average
effect?) and why they are close in these data.
\item Compare your estimates in (c) to the raw gap in (a). Two sentences: what accounts for the
difference, and what would the coefficient on \texttt{black} have to be interpreted as for this
regression to measure discrimination? (Munnell et al.\ argue one side; think about what an
omitted variable would need to look like.)
\item What share of the LPM's fitted probabilities lie outside $[0, 1]$? Which applicants are they?
\item Classification. Using the logit's predicted probabilities, build the confusion matrix at a
cutoff of $0.5$ and at $0.2$. Report accuracy, sensitivity (share of denials caught), and
specificity. Note that always predicting ``approve'' has accuracy $0.88$. One paragraph: which
cutoff would a lender's compliance officer prefer, which would a loan officer prefer, and why is
``accuracy'' the wrong summary for a rare outcome?
\end{enumerate}

\subsection*{2 Interest rates you only see for approved loans}

Return to the loan officers of Assignment 4, but now the outcome is the \emph{interest rate} on the
loan, which exists only if the loan is approved. $4{,}000$ applications, each randomly assigned to
one of $40$ officers with leniency $\ell_{j} \sim \mathcal{N}(0, 0.6^{2})$:
\begin{align*}
\text{approve}_{i} &= \mathbf{1}\{\,0.5\,\text{score}_{i} + \ell_{j(i)} + u_{i} > 0\,\} \\
\text{rate}_{i} &= 6 - 0.8\,\text{score}_{i} + 0.5\,e_{i}, \qquad \text{observed only if approved,}
\end{align*}
with $\text{score}_{i} \sim \mathcal{N}(0,1)$ and $(u_{i}, e_{i})$ standard normal with correlation
$\rho = 0.6$: applicants the officer is inclined to approve for reasons you do not see are also
charged more. The question is the effect of the credit score on the rate offered, $-0.8$.

\begin{enumerate}[label=(\alph*)]
\item Simulate. Regress rate on score using approved loans only. Sign and explain the bias before
you look at the number (Session 12: what is $\E[e \mid \text{approved}, \text{score}]$, and how does it
move with score?).
\item Heckman two-step: probit of approval on score and officer leniency; compute the inverse
Mills ratio $\lambda(\hat z) = \phi(\hat z)/\Phi(\hat z)$ from the probit index; regress rate on
score and $\hat\lambda$ in the approved sample. Compare the score coefficient to $-0.8$ and the
coefficient on $\hat\lambda$ to $\rho\sigma_{e} = 0.3$.
\item Redo (b) with a probit on score \emph{alone}, no leniency. Report the coefficient on score and
its standard error, and the correlation between $\hat\lambda$ and score in the approved sample.
One paragraph: what is the \emph{exclusion restriction} in this problem, why is officer leniency a
good one, and what happens to the two-step estimator without it?
\item Write the full-information likelihood for one observation (a probability for rejected
applications; a density times a conditional probability for approved ones) and maximize it with
\texttt{optim} (or \texttt{sampleSelection::selection} in R). Compare to (b).
\item Managerial summary, two sentences: a fintech competitor proposes to price loans using the
OLS from (a), estimated on the incumbent's approved book. What will go wrong when they lend to
applicants the incumbent would have rejected?
\end{enumerate}

\end{document}
